% Part: normal-modal-logic
% Chapter: axioms-systems
% Section: duals

\documentclass[../../../include/open-logic-section]{subfiles}

\begin{document}

\olfileid[zh]{nml}{prf}{dua}

\olsection{对偶\usetoken{P}{formula}}

\begin{defn}\ollabel{def:duals}
  \Ax{T}、\Ax{B}、\Ax{4} 和 \Ax{5} 的每个!!{formula}都有一个
  \emph{对偶}，用带下标的一个菱形表示，如下：
  \begin{align}
    \tag{\Ax{T_\Diamond}} p & \lif \Diamond p\\
    \tag{\Ax{B_\Diamond}} \Diamond\Box p & \lif p\\
    \tag{\Ax{4_\Diamond}} \Diamond\Diamond p & \lif \Diamond p\\
    \tag{\Ax{5_\Diamond}} \Diamond\Box p & \lif \Box p
    \end{align}
\end{defn}

上面的每一个对偶!!{formula}都由相应的!!{formula}将 $p$ 替换为
$\lnot p$、取逆否、把 $\lnot\Box\lnot$ 换成 $\Diamond$、再把
$\lnot\Diamond\lnot$ 换成 $\Box$ 而得到。\Ax{D}，即
$\Box!A \lif \Diamond!A$，在那个意义上是它自己的对偶。

\begin{prop}\ollabel{prop:dualsys}
  对 \olref[nml][prf][dua]{def:duals} 中的每个!!{formula}~$!A$：$\Log{K}!A = \Log{K}!A_{\Diamond}$。
\end{prop}
\begin{proof}
  习题。
\end{proof}
\begin{prob}
  证明 \olref[nml][prf][dua]{prop:dualsys}。
\end{prob}

\end{document}
